Un joueur participe à un jeu. Il dispose d'une pièce de monnaie équilibrée et de deux urnes A et B qui contiennent chacune des jetons dorés et des jetons argentés.

Il doit commencer par lancer la pièce de monnaie.

\begin{itemize}[left=5mm]
	\item S'il obtient le côté pile, il pioche un jeton dans l'urne A. Il a alors 60\,\% de chances de tirer un jeton doré.
	\item S'il obtient le côté face, il pioche un jeton dans l'urne B. Il a alors 80\,\% de chances de tirer un jeton argenté.
\end{itemize}

On considère l'expérience aléatoire qui consiste à lancer la pièce de monnaie puis à tirer un jeton dans une des deux urnes.

On considère les évènements suivants :

\begin{itemize}[left=15mm]
		\item A : le joueur pioche dans l'urne A;
		\item B : le joueur pioche dans l'urne B;
		\item D : le jeton obtenu est doré;
		\item $\overline{\mathrm{D}}$ : le jeton obtenu est argenté.
\end{itemize}

\medskip

\begin{enumerate}
	\item Construire un arbre pondéré traduisant la situation.
	\item Calculer $P\big(A \cap D\big)$, puis interpréter le résultat obtenu.
	\item Montrer que $P(D)=0,4$.
	\item Le joueur a obtenu un jeton doré.

	Quelle est la probabilité qu'il l'ait pioché dans l'urne B ?
\end{enumerate}

\needspace{10\baselineskip}
